\chapter{Logic Synthesis, Restructuring, and DFT}
\label{ch:logic}

\glance{OpenROAD can elaborate and map SystemVerilog itself (\pkg{src/syn}),
locally resynthesize critical cones via ABC (\pkg{src/rmp} + \pkg{src/cut}),
and insert scan chains (\pkg{src/dft}). The embedded ABC at \file{src/abc} is
\emph{out of scope} here --- see the separate ABC reference; this chapter only
covers how OpenROAD hands logic to ABC and stitches the result back.}

\section{Integrated synthesis: \pkg{src/syn}}

\id{syn::Synthesis} elaborates SystemVerilog (via slang) and maps to gates:

\begin{lstlisting}[style=ortcl]
sv_elaborate top.sv -top top_module   ;# slang options accepted
synthesize                            ;# map to stdcells -> odb
\end{lstlisting}

Limitations (per the module README): no latch mapping, always flattens, and
macro instance names may be lost. For a full RTL-to-GDSII flow most users
still rely on Yosys in ORFS (Ch.~\ref{ch:gui}); \pkg{src/syn} is the in-binary
path.

\section{Logic cuts: \pkg{src/cut}}

\pkg{src/cut} has \textbf{no user commands}. It is a library used by
\pkg{src/rmp}: it extracts a cloud of logic via the OpenSTA timing engine and
passes it to ABC (BLIF or a direct network), then reads the result back into
odb (including constants via tie cells).

\section{Restructure via ABC: \pkg{src/rmp}}

\id{rmp::Restructure} is the bridge to Berkeley ABC for \emph{local}
resynthesis:

\begin{enumerate}
  \item Select a timing-critical or area region (slack / depth thresholds).
  \item Extract it through \pkg{src/cut} and hand it to the embedded ABC.
  \item Run area or delay recipes; pick the best result.
  \item Stitch the mapped netlist back into odb (tie cells for constants).
\end{enumerate}

\begin{lstlisting}[style=ortcl]
restructure -liberty_file tech.lib -target delay \
            -slack_threshold 0 -depth_threshold 16 \
            -tielo_port TIELO/Y -tiehi_port TIEHI/Y
\end{lstlisting}

\cmd{-target area|delay} chooses the objective. ABC's own rewrite/mapping
algorithms are documented in the ABC manual, not here.

\section{Design for test: \pkg{src/dft}}

\id{dft::Dft} inserts scan infrastructure:

\begin{lstlisting}[style=ortcl]
set_dft_config -max_length 100 -clock_mixing no_mix
scan_replace          ;# swap flops for scan-equivalent cells
preview_dft           ;# dry-run the chain plan
insert_dft            ;# stitch scan chains, add SE/SI/SO ports
\end{lstlisting}

Elements added: scan-in / scan-out / scan-enable ports, scan cells, and one or
more shift-register chains. \cmd{-clock\_mixing} controls whether flops of
different clocks may share a chain.

\section{Summary}

\begin{table}[h]
\centering
\small
\begin{tabularx}{\linewidth}{@{}lllX@{}}
\toprule
\textbf{Module} & \textbf{Command} & \textbf{Purpose} & \textbf{External} \\
\midrule
\pkg{syn} & \cmd{sv\_elaborate}, \cmd{synthesize} & SV $\rightarrow$ gates in odb & slang \\
\pkg{cut} & (none) & extract logic cones for ABC & --- \\
\pkg{rmp} & \cmd{restructure} & local ABC resynthesis & ABC \\
\pkg{dft} & \cmd{scan\_replace}, \cmd{insert\_dft} & scan insertion & --- \\
\pkg{cgt} & \cmd{clock\_gating} & ICG insertion (Ch.~\ref{ch:cts}) & ABC proof \\
\bottomrule
\end{tabularx}
\end{table}
